\documentclass[12pt,a4paper,oneside]{report}

% Backend contribution report. Figure files used: logo.png,
% 02_scraping_workflow.png, 03_request_budget.png,
% 04_arrival_lag.png and 05_example_series.png.
% The original report is retained in the adjacent before_backend_revision file.

% -------------------------------------------------------------------
% Editable title-page details
% -------------------------------------------------------------------
\newcommand{\studentname}{Tanish Bansal}
\newcommand{\auravname}{Aurav Pratap Singh}
\newcommand{\rollnumber}{231075}
\newcommand{\supervisor}{Satya Prakash Singh}
\newcommand{\projectcourse}{Undergraduate Project Report (MTH496)}
\newcommand{\department}{Department of Mathematics and Statistics}
\newcommand{\institution}{Indian Institute of Technology, Kanpur}
\newcommand{\projecttitle}{Automated Web Scraping and Backend Data Infrastructure\\ for Agricultural Mandi Prices\\ in Uttar Pradesh and Rajasthan}

\usepackage[a4paper,margin=1in]{geometry}
\usepackage{amsmath,amssymb,amsfonts,amsthm,mathtools}
\usepackage{graphicx}
\usepackage{booktabs}
\usepackage{array}
\usepackage{longtable}
\usepackage{enumitem}
\usepackage{float}
\usepackage{caption}
\usepackage{subcaption}
\usepackage{hyperref}
\usepackage{microtype}
\usepackage{fancyhdr}
\usepackage{listings}
\usepackage{xcolor}
\usepackage{url}
\usepackage{setspace}
\usepackage{titlesec}
\graphicspath{{./}{figures/}}

\hypersetup{hidelinks,bookmarksnumbered=true,pdfauthor={\studentname},pdftitle={Automated Web Scraping and Backend Data Infrastructure for Agricultural Mandi Prices}}
\setstretch{1.10}\setlength{\headheight}{14pt}
\setlength{\parskip}{0.35em}
\setlength{\parindent}{1.5em}
\setlist[itemize]{leftmargin=2em,itemsep=0.2em,topsep=0.25em}
\setlist[enumerate]{leftmargin=2em,itemsep=0.25em,topsep=0.25em}
\captionsetup{font=small,labelfont=bf}
\numberwithin{equation}{chapter}
\numberwithin{figure}{chapter}
\numberwithin{table}{chapter}

\pagestyle{fancy}
\fancyhf{}
\fancyhead[L]{\small Mandi Price Backend Infrastructure}
\fancyhead[R]{\small Undergraduate Project Report}
\fancyfoot[C]{\thepage}
\renewcommand{\headrulewidth}{0.4pt}

\theoremstyle{definition}
\newtheorem{definition}{Definition}[chapter]
\newtheorem{example}{Example}[chapter]
\newtheorem{remark}{Remark}[chapter]

\DeclareMathOperator{\Var}{Var}
\DeclareMathOperator{\E}{E}
\DeclareMathOperator{\argmin}{arg\,min}
\DeclareMathOperator{\med}{median}
\newcommand{\R}{\mathbb{R}}
\newcommand{\hatt}[1]{\widehat{#1}}
\newcommand{\Rs}{\text{Rs.}}
\newcommand{\trans}{^{\mathsf T}}

% Figure that falls back to a placeholder box when the image is not uploaded yet.
\newcommand{\safefig}[2]{\IfFileExists{#1}{\includegraphics[width=#2\textwidth]{#1}}{%
  \fbox{\parbox[c][4.5cm][c]{0.8\textwidth}{\centering\small Figure file not found:\\\texttt{\detokenize{#1}}}}}}

\lstset{
  basicstyle=\ttfamily\small,
  columns=fullflexible,
  keepspaces=true,
  breaklines=true,
  frame=single,
  rulecolor=\color{black!30},
  showstringspaces=false,
  keywordstyle=\bfseries,
  commentstyle=\itshape\color{black!65}
}

\begin{document}

% =============================== TITLE ===============================
\begin{titlepage}
\thispagestyle{empty}

\centering
\IfFileExists{logo.png}{%
  \includegraphics[width=3.2cm]{logo.png}\par
}{%
  \framebox{\parbox[c][2.8cm][c]{2.8cm}{\centering\small\textbf{IIT Kanpur}\\\textbf{Logo}\\(\texttt{logo.png})}}\par
}
\vspace*{1.0cm}
{\LARGE\bfseries \projecttitle\par}
\vspace{1.8cm}
{\Large \studentname\par}
\vspace{0.25cm}
{\large Roll Number: \rollnumber\par}
\vspace{1.0cm}
{\large Project Supervisor: \supervisor\par}
\vspace{1.9cm}
{\large\bfseries \projectcourse\par}
\vspace{0.7cm}
{\large \department\par}
{\large \institution\par}

\vfill
{\large \today\par}
\vspace*{0.8cm}
\end{titlepage}

% ============================ FRONT MATTER ============================
\pagenumbering{roman}
\setcounter{page}{1}
\tableofcontents
\clearpage
\listoffigures
\clearpage
\listoftables
\clearpage

\chapter*{Abstract}
\addcontentsline{toc}{chapter}{Abstract}
Reliable market-price information requires more than a forecasting model: the underlying observations must first be collected, checked, stored and kept current. The project on agricultural mandi prices is organised into several phases: source investigation and collection, data validation and storage, scheduled processing and backend maintenance, forecasting, and website development. I handled the collection pipeline, data validation and provenance, SQLite storage, persistent jobs, scheduling, restart recovery, data-serving support and routine operation. \auravname\ handled forecasting methodology and model selection, website design and front-end development. This report documents the phases handled by me; the forecasting and website phases are mentioned only where needed to explain how they use the backend.

The collection work begins with measured properties of the source: reports accept all districts of a state, one commodity at a time, a year-long date interval and 5,000-row pages; access tickets are tied to the report filters; and requests are throttled after an initial burst. These constraints motivate one state--commodity report per year for initial collection, followed by recent-date updates. The pipeline validates each row, converts compatible price units, preserves missing dates, quarantines unusable records and records corrections without duplicating observations. Jobs and checkpoints are stored alongside the data so that failed work can be resumed.

An in-process scheduler queues daily updates at 07:00 Indian Standard Time, a wider census on Sunday at 03:00, and a local downstream-result rollover at 00:05. A single worker executes eligible jobs, backs off on throttling, defers failed commodities and recovers interrupted runs. The report explains the distinction between a scheduled trigger, successful collection and the freshness of the source's own reports. Historical collection evidence and successful runs recorded on 6 November 2026 demonstrate the backend in operation. The system currently runs on one personal computer; an always-on public deployment remains future work.

\clearpage
\chapter*{Declaration and Contribution Statement}
\addcontentsline{toc}{chapter}{Declaration and Contribution Statement}
I hereby declare that this report documents the phases of the agricultural mandi-price project handled by me: backend data acquisition, validation, database management, scheduled processing, recovery and operational maintenance. The project phases and their respective responsibilities are described in Chapter~1.

The forecasting and website-development phases were handled by \auravname. This report describes my backend work and the data interfaces connecting it to those phases. Published ideas and methods are acknowledged through citations and the bibliography. Numerical collection and operational observations are associated with the dated evidence described in this report.

\vspace{1.2cm}
\noindent Name: \studentname \hfill Roll Number: \rollnumber
\vspace{1.0cm}

\noindent Signature: \rule{7cm}{0.4pt}
\vspace{0.8cm}

\noindent Date: \rule{4cm}{0.4pt}

\clearpage
\chapter*{Acknowledgements}
\addcontentsline{toc}{chapter}{Acknowledgements}
I thank my project supervisor, \supervisor, and the \department, \institution, for guidance. I thank \auravname\ for handling the forecasting and website-development phases and for discussions on connecting those phases to the backend data pipeline.

\clearpage
\pagenumbering{arabic}
\setcounter{page}{1}

\chapter{Problem, Project Phases and Report Scope}

\section{The problem}
A farmer deciding where and when to sell a crop needs reliable prices for nearby mandis. Agmarknet publishes reported minimum, maximum and modal prices, but obtaining a report requires a sequence of filter selections, a date interval and a CAPTCHA \cite{agmarknet}. Repeating that process manually across many commodities and markets is slow. A forecasting website also needs these observations to arrive in a consistent, queryable form; otherwise a working interface can display old, incomplete or incompatible prices.

The project therefore progresses through connected phases, from collecting source reports and maintaining a trustworthy dataset to generating forecasts and presenting them on a website. I handled the data pipeline and backend operation, while \auravname\ handled forecasting and website development. This report covers the phases handled by me. The backend problem is to collect the source's reports efficiently, account for every received row, survive interruptions and make the saved observations available without making a visitor wait for source downloads.

\section{Project phases and responsibilities}
Table~\ref{tab:division} lists the main phases and who handled them. The technical chapters, evaluation and conclusion of this report focus on my phases.

\begin{table}[H]
\centering
\caption{Project phases, responsibilities and scope of this report}
\label{tab:division}
\small
\begin{tabular}{@{}p{0.23\textwidth}p{0.18\textwidth}p{0.50\textwidth}@{}}
\toprule
Phase & Handled by & Scope \\
\midrule
Source investigation and collection & \studentname & Source API behaviour, report requests, CAPTCHA and tickets, pagination, throttling and resumable scraping. \\
Data validation and storage & \studentname & Normalisation, unit handling, provenance, SQLite storage, quarantine and correction history. \\
Scheduled processing and backend operation & \studentname & Persistent jobs, scheduling, retries, restart recovery, data-serving support, freshness metadata and operational maintenance. \\
Forecasting & \auravname & Forecasting mathematics, model choice and comparison, and predictive evaluation. \\
Website development & \auravname & Website design, front-end implementation, controls, charts and graphical presentation. \\
\bottomrule
\end{tabular}
\end{table}

\section{Objectives of my contribution}
\begin{enumerate}
  \item establish the actual behaviour of the source API, including pagination, access tickets, empty reports and throttling;
  \item reduce the number of report requests while preserving the source's market, commodity, variety and grade distinctions;
  \item validate, normalise and store observations with an audit trail;
  \item run daily and weekly collection through persistent, resumable jobs;
  \item preserve correct behaviour after failures, restarts and repeated requests;
  \item support the website with stored data, freshness information and responsive backend reads;
  \item evaluate collection effort and operational reliability, rather than predictive accuracy.
\end{enumerate}

\section{Scope and terminology}
The maintained application covers Uttar Pradesh and Rajasthan. A small earlier Maharashtra pilot remains in storage but is outside the active two-state service. The broader command-line collection tools and the application-specific job service share storage facilities; the live application does not drain the old nationwide planning queue.

A \emph{window} is a state, commodity and date interval fetched as a report. A \emph{series} is a fixed state--district--market--commodity--variety--grade combination. A \emph{job} is a persisted unit of application work, such as a state's daily update. These have different lifetimes: a job may contain many windows, and each window may update many series.

Forecasting formulas, model comparisons, interval calibration and front-end design are omitted from this report. Downstream result storage is discussed only to explain how the backend avoids serving results for an obsolete data revision. The documented service is the local Python application in \path{agmarknet_collector}; the separate \path{harvest} web project is not the daily ingestion service examined here.

\section{Report organisation}
Chapter~2 explains source access and the scraping design. Chapter~3 develops validation and database management. Chapter~4 details scheduling, queue selection, retries and recovery. Chapter~5 explains backend integration and maintenance of the running service. Chapter~6 examines reporting lag and data coverage. Chapter~7 presents collection evidence and verification. Chapter~8 gives limitations and future work, followed by the conclusion. The appendices provide database, API, algorithm and reproduction details.

\begin{figure}[H]
\centering
\fbox{\begin{minipage}{0.89\textwidth}
\centering
\textbf{Backend contribution documented here}\\[0.4em]
Agmarknet reports $\longrightarrow$ source client and validation\\[0.3em]
$\longrightarrow$ SQLite observations, provenance and job checkpoints\\[0.3em]
$\longrightarrow$ scheduled refreshes and backend data services\\[0.8em]
\hrule
\vspace{0.7em}
\textbf{Forecasting and website phases: \auravname}\\[0.4em]
Forecasting methodology and model selection; website and front end\\[0.3em]
consume the backend's verified observations and freshness metadata.
\end{minipage}}
\caption{Data flow between the backend phases documented here and the forecasting and website phases.}
\label{fig:boundary}
\end{figure}

\chapter{Source Access and Automated Web Scraping}

This chapter describes the source investigation and collection design. API behaviour and throughput figures are observations from the project probes, not guarantees about the portal. The collection measurements retained from the original experiment refer to the 27 October 2026 snapshot unless a later date is stated.

\section{The report interface}
The portal's ``Daily Price and Arrival Report'' page is a single-page application that calls a JSON API. Calling the structured endpoints avoids extracting table contents from rendered HTML, an approach consistent with API-based web scraping \cite{glez2014}. A report is requested with a set of filters: price/arrival type, commodity group and commodity, variety, grade, state, district, market, and a date range. The collector uses two kinds of calls:
\begin{enumerate}
  \item a \emph{catalogue} call returning the options for every filter (states, districts, markets, commodities, groups, varieties, grades) and the allowed date range;
  \item the \emph{report} call, which returns a page of rows together with pagination information (total count, total pages, current page).
\end{enumerate}
Each row of a report holds the arrival date, state, district, market, commodity, variety, grade, minimum, maximum and modal price, the price unit, and (for combined reports) the arrival quantity.

\section{CAPTCHA and the access ticket}
Before a report is served, the client must obtain a CAPTCHA image, answer it, and receive a short-lived \emph{ticket}. Three observations from experiments shaped the design:
\begin{itemize}
  \item the ticket is \emph{bound to the exact filters}: if the filters change, the portal answers ``Report filters changed. Please solve captcha again.'' Hence one CAPTCHA is needed for every distinct report;
  \item once a ticket is issued for a report, it can be reused for all pages of \emph{that} report;
  \item a wrong answer is reported as an invalid-CAPTCHA error, with an error code that is shared with throttling (Section~\ref{sec:throttle}).
\end{itemize}
The images contain six plain characters. They are solved locally by a glyph-template matcher: the font glyphs are rendered at several sizes, the image is segmented into six characters and each is matched against the templates (Pillow and NumPy). The solver rejects images it cannot segment instead of guessing, uses no remote solving service, and is limited to three fresh challenges per report attempt. Section~\ref{sec:perf} reports the diagnostic events observed. Figure~\ref{fig:scrapeflow} summarises the whole flow.

\begin{figure}[H]
\centering
\safefig{02_scraping_workflow.png}{0.98}
\caption{Flow of one report request: obtain a CAPTCHA, solve it locally, receive a ticket, then fetch the pages of the report with that ticket; errors are classified as CAPTCHA, throttling, outage or format change.}
\label{fig:scrapeflow}
\end{figure}

\section{Measured behaviour of the report API}
Controlled probes (a handful of requests each) established the facts in Table~\ref{tab:apifacts}, which were then used as design constraints.

\begin{table}[H]
\centering
\caption{Measured behaviour of the report API and the design consequence}
\label{tab:apifacts}
\small
\begin{tabular}{@{}p{0.46\textwidth}p{0.46\textwidth}@{}}
\toprule
Finding & Consequence \\
\midrule
A report for \emph{all districts} of a state and one commodity can be requested in one call. & One report per (state, commodity), not per market. \\
The \texttt{commodity} parameter accepts exactly one commodity id. & Commodities cannot be batched in one report. \\
A date span of up to 365 days is served in one report; any year since 2021 can be requested. & A year of history is one report per commodity. \\
Pages of up to 5,000 rows are accepted. & Uttar Pradesh wheat ($\approx 85{,}000$ rows per year) needs 17--18 pages and one CAPTCHA. \\
A commodity with no records returns HTTP 404 with the text ``No data available''. & An empty report is a valid result, not an error. \\
An ``all states'' option exists in the state filter. & An all-states design is possible (Chapter~8). \\
\bottomrule
\end{tabular}
\end{table}

\section{Throttling and its classification}\label{sec:throttle}
During an early probe about one third of requests seemed to fail the CAPTCHA. The cause turned out to be a different condition: after a burst of roughly 100 requests the portal answers with HTTP status 400, error code \texttt{TOKEN\_OR\_CAPTCHA\_REQUIRED} and the message ``too many requests'', and then sustains only about 1--2.5 requests per minute. The code is the same as for a wrong CAPTCHA, so the two must be told apart by the message text. Treating throttling as a CAPTCHA failure wastes CAPTCHA attempts and makes the situation worse. The client therefore distinguishes four conditions and handles them differently (Table~\ref{tab:errors}).

\begin{table}[H]
\centering
\caption{Error classes and the collector's response}
\label{tab:errors}
\small
\begin{tabular}{@{}p{0.2\textwidth}p{0.36\textwidth}p{0.36\textwidth}@{}}
\toprule
Condition & Detection & Response \\
\midrule
Wrong CAPTCHA & invalid-captcha message & Fetch a fresh image; give up on that report after three. The report is \emph{deferred}. \\
Throttling & ``too many requests'' message & Stop the sweep; wait 15 minutes, then 30; slow the request pace. \\
Outage & Network error or HTTP 5xx after five attempts & Wait and retry later; do not change the pace. \\
Format change or denial & Non-JSON or unexpected shape; access denied & Pause collection and show it on the admin page. \\
\bottomrule
\end{tabular}
\end{table}

\section{The request budget: why design matters}
The cost of a plan is the number of CAPTCHA-solved reports (plus extra pages). The legacy plan (one report for each state--commodity--month) is
\[
  36\ \text{states}\times 605\ \text{commodities}\times 13\ \text{months}=283{,}140\ \text{reports}.
\]
At roughly one to two requests per minute after the burst, with normally more than one HTTP exchange per report, this is months of continuous running. The plan used here collects, for a chosen state, \emph{one report per commodity for the whole year}. Uttar Pradesh needed 408 commodity windows (commodities in groups likely to report prices) and Rajasthan 150 (Table~\ref{tab:budget}). Of the roughly 1,100 report windows completed for the two states, 478 (43\%) were empty, which is why a state added later starts from the commodities already seen elsewhere (Section~\ref{sec:jobs}). Figure~\ref{fig:budget} compares the plans. Its all-states value is a historical planning estimate, not a measured run or the count of distinct commodity identifiers; a nationwide design must account separately for page requests and any date splitting.

\begin{table}[H]
\centering
\caption{Report windows needed by different plans}
\label{tab:budget}
\small
\begin{tabular}{@{}p{0.5\textwidth}rp{0.3\textwidth}@{}}
\toprule
Plan & Reports & Note \\
\midrule
All states $\times$ commodities $\times$ months & 283,140 & Not feasible \\
Uttar Pradesh, one report per commodity per year & 408 & Used; about 12 hours \\
Rajasthan, seeded commodity list & 150 & Used; about 2.5 hours \\
All states, one report per commodity (estimate) & 1,500--2,000 & Proposed (Chapter~8) \\
\bottomrule
\end{tabular}
\end{table}

\begin{figure}[H]
\centering
\safefig{03_request_budget.png}{0.85}
\caption{Report windows required by each collection plan (logarithmic scale).}
\label{fig:budget}
\end{figure}

\section{Collector design}
\subsection{One window}
A window is downloaded page by page. The application requests up to 5,000 rows per page and checks the page number, total row count and total page count throughout the download. A repeated non-empty page, an unexpected empty page or inconsistent totals cause the window to fail rather than accepting a truncated report. A report above 250,000 rows is divided into two date intervals recursively; a single-day report above that limit requires attention.

An exact window key contains the commodity and interval within its job. If that key is already complete, it is skipped. After a full, consistent download, validation and observation storage are transactional. The application records the window checkpoint after that storage transaction, so a crash in between can cause a safe re-fetch; the idempotent observation keys prevent duplication. This is an at-least-once execution design, not a claim of exactly-once network access.

\begin{lstlisting}[caption={Fetching and checkpointing a report window}]
function fetch_window(state, commodity, start, end, job):
    if exact_window_complete(job, commodity, start, end): return
    create pending batch and window records
    rows = []; page = 1
    loop:
        body = client.report(filters, page, limit=5000)
        verify page number and stable pagination totals
        if first page reports more than 250000 rows:
            split date interval; fetch both halves; checkpoint parent
            return
        reject repeated pages or premature empty pages
        rows += body.rows
        if page >= body.total_pages: break
        page += 1
    verify len(rows) == reported_total
    transaction: validate rows; upsert good rows; quarantine bad rows;
                 record revisions and batch outcome
    mark this exact application window complete
    invalidate backend list caches
\end{lstlisting}

\subsection{A sweep with deferral}
A \emph{sweep} fetches one window per commodity. A CAPTCHA failure defers only that commodity; throttling stops the sweep and defers everything that remains; an outage does the same without changing the pace.

\begin{lstlisting}[caption={Sweep over commodities with deferral and pace adaptation}]
function sweep(state, commodities, start, end, job):
    deferred = []; throttled = False; outage = False
    for each commodity c in the selected list:
        try: fetch_window(state, c, start, end, job)
        except CaptchaFailed:  deferred.append(c)
        except Throttled:      deferred = commodities[c:]; throttled = True; break
        except Unavailable:    deferred = commodities[c:]; outage = True; break
    if not outage:
        pace = pace * 1.5 if throttled else pace * 0.9      # clamp to [4, 30] seconds
    return deferred

# throttling/outage: usually 15 then 30 minutes
# CAPTCHA-only deferral: normally 10 then 20 minutes
# completed identical windows are skipped on retry
\end{lstlisting}

\subsection{Pace learning}
The request interval starts at 7.5 seconds, is multiplied by 1.5 after a throttled sweep and by 0.9 after a clean one, and is kept between 4 and 30 seconds. It is stored in the database, so a restart continues with the learned pace. This is a simple multiplicative-increase / multiplicative-decrease rule; the portal's real limit is not published.

\subsection{Single worker and resumability}
A single worker thread runs one job at a time, and an operating-system file lock prevents two collector processes from running together. Jobs (kind, state, status, attempts, retry time) and completed windows are stored in the database. After a crash or restart, a job left ``running'' is re-queued, and its already completed windows are not fetched again. Eligible non-history jobs are selected before queued history builds. An already running history build is not interrupted; the queue policy is non-preemptive. Chapter~4 gives the exact selection and recovery rules.

\section{Responsible use}\label{sec:responsible}
The portal serves public government data. The collector keeps to one request at a time, learns a conservative pace, backs off when told to, does not use third-party CAPTCHA-solving services, and archives raw pages only for the history build to avoid repeated downloads. Before any wider deployment the portal's terms of use should be checked and the request rate kept within what the operators tolerate.

% ============================== CHAPTER 4 ==============================
\chapter{Data Validation and Database Management}

\section{Why the website needs persistent storage}
Downloading a source report for every visitor would repeat expensive requests, consume the source's allowance and make response times depend on CAPTCHA recognition and network availability. The backend therefore separates collection from reading. Collection writes verified observations to \path{output/agmarknet.sqlite}; the running application reads those observations even while a later update is queued or waiting for retry. A source outage can stop new ingestion without deleting the saved history.

SQLite is sufficient for the current single-machine application because both collection and data-serving run near the same database file. Write-ahead logging (WAL) lets readers generally continue while a writer commits changes \cite{sqlitewal}. It does not create multiple simultaneous writers: collection remains serial and short transactions limit contention. Each store connection has a 30-second database timeout and is closed after use.

\section{Row normalisation and validation}
The source contains date strings, numbers with commas, inconsistent whitespace and price-unit labels. Storage converts accepted records into a consistent representation before exposing them downstream:
\begin{enumerate}
  \item Parse supported date formats into ISO dates and require the date to lie within the requested interval.
  \item Collapse repeated whitespace in dimension fields. Require state, district, market and commodity; replace a missing variety or grade with ``Not reported''.
  \item Verify that state and commodity match the requested catalogue entries, and verify market membership when a market filter was requested.
  \item Parse numeric prices after removing commas. The modal price is required; minimum and maximum can be absent. Reject negative and non-finite values while retaining a genuine reported zero.
  \item Convert recognised compatible units to INR per quintal. A price per kilogram is multiplied by 100; a price per tonne is multiplied by 0.1. A bundle price is not converted without a defensible weight.
  \item Where minimum and maximum exist, require minimum $\le$ modal $\le$ maximum.
\end{enumerate}

The source spells its modal-price field \texttt{model\_price}; the normaliser maps this exact source key to \texttt{modal\_price}. Treating a guessed spelling as equivalent would silently lose valid data. Likewise, HTTP 404 is accepted as an empty report only for the verified ``No data available'' response, not for a missing route or malformed body.

\section{Quarantine instead of silent deletion}
A rejected row is retained with its raw contents, batch identifier, reason and creation time in \texttt{quarantine}. Valid rows from the same report may still be stored, but the batch becomes \texttt{complete\_with\_quarantine}. This distinguishes accounting for all downloaded rows from accepting all of them as usable prices. Duplicate natural keys inside one report are also recorded for review rather than silently allowing the later quotation to overwrite the earlier one.

In the historical two-state snapshot, 3,414 rows were quarantined: 3,387 had the unsupported ``Rs./Bundle'' unit and 27 violated price ordering. These are collection-quality results, not a claim that a later database will contain the same counts. Large movements in otherwise valid prices are retained; review flags in the export distinguish an unusual observation from a demonstrably invalid one.

\section{Observation keys and series identities}
An observation's natural key is the ordered tuple
\[
(\text{date},\text{state},\text{district},\text{market},
 \text{commodity},\text{variety},\text{grade}).
\]
After whitespace normalisation, text values are case-folded for key construction. The key is stored as a JSON-encoded tuple and enforced as the primary key of \texttt{observations}. Re-fetching an overlapping date interval therefore updates the same observation instead of inserting a second copy.

A series uses the same dimensions without the date. Its identifier is a SHA-256 hash of the JSON-encoded, case-folded six-field tuple; the readable dimensions are stored separately in \texttt{series\_state}. Varieties and grades remain separate, so the backend does not average unlike quotations into a new target. Any statistical aggregation is a downstream modelling decision and is outside my contribution.

\section{Idempotent updates and correction history}
For an existing observation, the storage layer compares minimum, maximum, modal price and normalised unit with the saved record. An unchanged re-fetch refreshes the record's payload and provenance without increasing the series revision. A changed price or unit writes the previous and new payloads into \texttt{revisions}. A newly inserted or changed observation advances its series revision once per affected series in that window.

The revision counter is an integration signal: a downstream result generated for an older revision cannot be assumed current. This is distinct from the source date and retrieval time. The historical snapshot recorded no corrections to already stored prices; that observation does not remove the need to support corrections in subsequent operation.

\section{Transaction boundaries and consistency}
Within \texttt{Store.finish}, accepted observations, revision records, series revision counters, quarantine rows and batch totals are written in a database transaction. A database exception rolls the transaction back. Previously committed windows remain available if a later window fails.

The \texttt{app\_windows} completion marker is committed separately after \texttt{finish}. Consequently, a crash can leave valid data stored while its application checkpoint is still pending. Recovery may download that window again, and the natural-key upsert makes this safe. A job can therefore have useful committed data even when its overall status is \texttt{retry\_wait} or \texttt{paused}; overall success means all selected windows and the invoked downstream refresh have finished.

\section{Database responsibilities}
The application uses one SQLite database, with tables serving different responsibilities (Table~\ref{tab:tables}). The downstream cache is included to explain lifecycle management, not the forecasting algorithm.

\begin{table}[H]
\centering
\caption{Tables supporting ingestion and website operation}
\label{tab:tables}
\small
\begin{tabular}{@{}p{0.28\textwidth}p{0.64\textwidth}@{}}
\toprule
Table & Responsibility \\
\midrule
\texttt{observations} & Canonical dated prices with six market/product dimensions and source payload. \\
\texttt{series\_state} & Stable series identity, readable dimensions and data revision. \\
\texttt{batches}, \texttt{pages} & Source-request scope, outcome, counts and archived-page metadata. \\
\texttt{quarantine}, \texttt{revisions} & Rejected rows and previous values of corrected observations. \\
\texttt{events} & Collector batch events. \\
\texttt{app\_runs}, \texttt{app\_windows} & Persistent job queue, attempts, retry times and exact-window checkpoints. \\
\texttt{app\_settings} & Pause flag, request pace and per-state update/history markers. \\
\texttt{app\_diagnostics} & Request failures and diagnostic categories linked to jobs. \\
\texttt{app\_arrivals} & Per-date observation counts after successful state updates. \\
\texttt{app\_forecasts} & Results from the forecasting phase indexed by series, target date, data revision and model version. \\
\bottomrule
\end{tabular}
\end{table}

\section{Provenance, indexes and bounded archives}
Each accepted payload retains the source URL, request batch, retrieval timestamp, original dimension text and original price unit alongside the normalised fields. Initial history windows archive redacted raw JSON pages on disk, with relative path, SHA-256 checksum and counts in \texttt{pages}. CAPTCHA answers, keys and report-access tickets are removed from archived bodies. Routine updates preserve validated row provenance but do not archive every raw response, limiting repeated raw-page growth. This does not place a strict upper bound on total database size; old jobs, diagnostics and downstream caches still need a retention policy.

The composite observation index orders state, district, market, commodity, variety, grade and date, matching a history lookup for one series. The series primary key supports revision lookup, and the downstream result's composite key supports exact version selection. These structures support reads without re-parsing raw report files.

\section{Backup and restoration considerations}
The database and any WAL contents form the live storage state. Copying only the main SQLite file while the service is writing can omit recent commits; a backup should use SQLite's backup API or be taken after a clean stop \cite{sqlitebackup}. Raw archives and the filter catalogue should be retained separately when reproducible source evidence is needed. The current application does not implement an automated backup schedule or a tested remote restore service; those are deployment tasks rather than completed features.

\chapter{Scheduling, Job Selection and Recovery}

\section{Scheduler and worker are separate}
\label{sec:jobs}
The runtime owns an APScheduler \texttt{BackgroundScheduler} and one collector-worker thread. A scheduled callback only inserts or reuses a job in SQLite; it does not download reports itself. The worker drains eligible jobs and, when none remain, waits up to 10 seconds before checking again. A queued trigger time is therefore not a promised completion time.

The scheduler uses Indian Standard Time (UTC+05:30). Database timestamps are stored as UTC ISO strings; the administration interface displays them locally. Table~\ref{tab:jobs} lists the registered triggers. Initial history is requested on startup when a state lacks a completed build, rather than by a periodic cron trigger.

\begin{table}[H]
\centering
\caption{Registered application jobs and their purpose (IST)}
\label{tab:jobs}
\small
\begin{tabular}{@{}p{0.20\textwidth}p{0.22\textwidth}p{0.49\textwidth}@{}}
\toprule
Job & Trigger & Work queued \\
\midrule
Daily update & Every day 07:00 & One update per tracked state, re-fetching recent dates for active commodities. \\
Weekly census & Sunday 03:00 & One wider update per state, including additional commodities. \\
Local rollover & Every day 00:05 & One shared refresh of results from the forecasting phase for the new target date; no source scraping. \\
History build & Startup if needed; manual rebuild & Initial whole-period windows for one selected state. \\
\bottomrule
\end{tabular}
\end{table}

\section{Selecting the work inside a job}
For a daily update, active commodities are those with an observation in that state within the last 60 days. The requested interval is \emph{today minus 14 days through today}; both endpoints are included, so it can cover 15 distinct calendar dates. This overlap is deliberate because reports for older dates can arrive after earlier collection attempts.

The weekly census uses today minus 30 days through today and adds up to 40 commodities outside the active list. Extras are ordered by their last successful check, with never-checked commodities first and names breaking ties. This gradually discovers new or seasonal commodities without scanning the full catalogue every day. The filter catalogue is cached in \path{catalogue.json}; automatic detection of newly added catalogue entries is not implemented, so that catalogue needs explicit refresh when the source changes.

The initial history build defaults to 365 inclusive calendar dates. For Uttar Pradesh, priority crops are fetched first and groups unlikely to report prices are left to census checks. Rajasthan starts from commodities observed in other states plus a list of local specialities. Initial seeding reduces empty requests; it is not proof that every possible commodity is already covered.

\section{Queue deduplication and selection algorithm}
Enqueueing uses a \texttt{BEGIN IMMEDIATE} database transaction. If the same kind and state already has a \texttt{queued}, \texttt{running} or \texttt{retry\_wait} job, its identifier is returned. Otherwise a new job is inserted. Forecast-result refresh jobs have no state and are shared. An update and weekly census are different job kinds and may both be queued.

For execution, the worker selects jobs with status \texttt{queued}, or \texttt{retry\_wait} whose \texttt{retry\_at} has arrived. It orders eligible non-bootstrap jobs before bootstrap jobs and then by creation time. This is a non-preemptive, oldest-first queue within those two groups, not a CPU time-slicing algorithm. A running history build or long census can delay later work. There is no separate priority for local result refreshes among non-bootstrap jobs.

\begin{lstlisting}[caption={Selection of the next eligible backend job}]
acquire in-process mutex without waiting; otherwise return
acquire exclusive collector file lock; otherwise return

if collection is paused:
    allow only local downstream-refresh jobs
else:
    allow all job kinds

choose a queued job or due retry-wait job
order by (is_bootstrap, created_time)
if none exists: release locks; return
mark chosen job running
execute it while holding the collector locks
record success or the appropriate retry/pause transition
release locks
\end{lstlisting}

The mutex prevents concurrent work within the process. The file lock on \path{output/.collector.lock} prevents another cooperating collector process from using the same output directory at the same time. This is appropriate for the local filesystem deployment; it is not a distributed locking service.

\section{Job state transitions and checkpoints}
New work enters \texttt{queued}, becomes \texttt{running}, and ends in \texttt{succeeded}, \texttt{retry\_wait} or \texttt{paused}. A retry time makes deferred work eligible again; it does not busy-wait inside the download loop. Failure details, attempt count, timestamps and progress remain queryable.

Checkpoints are per exact commodity/date window, not per individual downloaded page. A retry skips an identical completed window, but restarts an incomplete one from its report pages. If a retry crosses midnight, the update recomputes its interval from the current date; changed interval keys can be fetched again. Natural-key upserts prevent duplicated observations in either case. This distinction explains what resumability guarantees and what it does not.

\section{Failure classification and retry delays}
The client performs bounded request-level retries for network failures and certain HTTP errors, normally up to five attempts with short exponential waits and consideration of \texttt{Retry-After}. Longer source cooldowns are handed back to the job service. There are therefore two layers: quick request recovery and persisted job deferral.

\begin{table}[H]
\centering
\caption{Persisted job response after request-level handling}
\label{tab:retry}
\small
\begin{tabular}{@{}p{0.23\textwidth}p{0.67\textwidth}@{}}
\toprule
Failure & Worker policy \\
\midrule
CAPTCHA exhaustion & Defer the commodity and continue the sweep. If it is the only deferral cause, retry after a base 10 minutes, then 20 minutes on subsequent failures. \\
Throttling & Stop the sweep and defer the remainder. Usually retry after 15 minutes, then 30 minutes; a source-specified delay can be longer. \\
Network/source outage & Stop and defer remaining work, usually 15 then 30 minutes, without changing the learned request pace. \\
Access denied or response-format change & Pause collection immediately, retaining saved data and diagnostics. \\
Other exception & Retry once after 30 minutes; pause if it fails again. \\
\bottomrule
\end{tabular}
\end{table}

For a deferred or rate-limited job, the delay is its supplied base delay multiplied by $2^{\min(k,1)}$, where $k$ is the number of previously recorded failures. The configured deferral limit is 12 automatic retries after failures before collection pauses. A sweep with mixed failures uses the delay returned by the condition that stopped it, so not every CAPTCHA-containing run has the CAPTCHA-only delay.

The collection pause is global: a fatal failure for one state prevents scheduled collection for both states. Local downstream-result refreshes remain eligible. A manual collection action clears the pause and queues or reuses work, but does not change all historical paused job records back to queued.

\section{Adaptive pacing}
The persisted request interval starts at 7.5 seconds. After a throttled sweep it is multiplied by 1.5, capped at 30 seconds; after a sweep without throttling or outage it is multiplied by 0.9, floored at 4 seconds. The adjusted interval is used when constructing the next job attempt's client. An outage is excluded from pace adaptation because it is not evidence of the source's sustainable request rate.

The new pace is not a replacement for cooldowns: multiplying the interval alone cannot recover an already exhausted request allowance. The policy combines conservative spacing with longer deferred waits, and never introduces concurrent report downloads.

\section{Restart recovery and catch-up}
At startup the application first attempts to acquire the collector file lock. If successful, jobs left \texttt{running} are returned to \texttt{queued}; their committed data and window checkpoints remain. The three cron triggers are then registered again, the scheduler starts, and startup catch-up considers each state:
\begin{enumerate}
  \item if history is incomplete, enqueue a bootstrap job;
  \item otherwise, if there is no successful-update timestamp or it is older than 24 hours, enqueue one update;
  \item enqueue one shared local downstream refresh.
\end{enumerate}

The work queue is persistent, but the cron trigger registration itself uses the scheduler's in-memory store. On restart, catch-up does not reconstruct every missed daily or weekly trigger. It queues one daily update when needed; a missed weekly census is not automatically replayed. This prevents a burst of redundant daily requests but leaves a wider discovery sweep to the next scheduled or manual census.

\section{Late triggers and duplicate callbacks}
The scheduler is configured with \texttt{coalesce=True}, \texttt{max\_instances=1} and \texttt{misfire\_grace\_time=3600} \cite{apscheduler}. Coalescing merges accumulated due executions of a scheduler callback; the grace period permits eligible late callback execution within one hour. The instance limit applies to that callback, which enqueues work, while the worker and locks enforce collection serialisation. These options do not wake a sleeping computer or keep a stopped process alive.

\section{Why one scheduled collection a day}
The initial design checked three times a day. In the historical probes a full recent-date Uttar Pradesh recheck changed 32 of 16,852 stored series (about 0.2\%), while a Rajasthan recheck two to three hours after its history build changed 264 of 4,554 (about 5.8\%). Those observations, together with reporting delays of days, motivated one daily trigger plus a wider weekly sweep. They do not establish that 07:00 is optimal for every state. The arrival log provides evidence for later timing changes.

\chapter{Backend Integration and Daily Service Operation}

\section{Serving saved data to the website}
The integrated local application uses FastAPI and uvicorn for its HTTP service \cite{fastapi}. My contribution at this boundary is the availability, correctness and freshness of the backend data. Website layout, controls, chart rendering and front-end implementation were handled by \auravname.

The history endpoint obtains a stable series identity, resolves its dimensions and returns observations in date order, optionally limited to a recent interval. A CSV route exports the same stored history with units and market/product dimensions. Selections and status routes expose the available series, tracked states, collection progress and last successful check. Pages therefore read from local storage instead of issuing new Agmarknet requests.

The farmer page polls the backend every 60 seconds and reloads selections when the last-successful-update marker changes. This is the consumption behaviour of the website interface, included here to explain how newly committed backend data becomes visible. A page request can queue a missing local result refresh, but it does not initiate source scraping.

\section{Downstream integration without a forecasting claim}
The job service calls the forecasting component after history and update sweeps, and at the midnight rollover. The statistical method was developed in the forecasting phase handled by \auravname. From the backend's perspective, it consumes date-ordered observations and produces a result payload.

Cached result identity includes series, target date, data revision and model version. Only an exact match is treated as current. Before saving a computed result, the service checks the series revision again; a changed revision invalidates the computation. Changes to downstream logic must also change its model-version identifier. Missing current results are represented as updating instead of being silently substituted with a result for another date.

This is lifecycle coordination and storage management. It does not evaluate which model is best or reproduce its mathematical derivation. A downstream refresh failure can also keep an ingestion job from being marked fully successful even after some or all observations have committed, so job detail and stored data must be inspected together.

\section{Keeping reads responsive during collection}
State selections and summary queries can scan many series. The service uses a short-lived in-memory cache with a 60-second lifetime, current-date check and generation counter. Saving a window invalidates the generation. In the live server, an expired entry can be returned while one background thread rebuilds it; repeated readers share that rebuild instead of each launching a new scan.

This deliberately permits a slightly older list during refresh while preserving the committed database. Collection-dependent processing uses a fresh-read context, so it does not build downstream outputs from the temporarily stale display cache. Startup warms the farmer lists, and successful jobs warm them again. History reads use the database directly. The report's historical 10--30 millisecond response observation refers to warmed list/API reads on the local machine, not a guaranteed latency for every route or deployment.

\section{Freshness is more than the last successful run}
The backend distinguishes three times:
\begin{enumerate}
  \item \textbf{Observation date:} the market date to which a reported price belongs.
  \item \textbf{Retrieval time:} when that source row was fetched and validated.
  \item \textbf{Last successful update:} when a state's complete update workflow finished.
\end{enumerate}
A successful run today can retrieve a report whose newest usable observation is several days old. Conversely, a retrying run can already have committed useful new prices before it finishes. The public status reports overall completion, while provenance and job progress explain partial progress. The latest date alone also cannot establish that all commodities or markets have reported on that date.

\section{Administration and monitoring support}
The backend exposes job status, next scheduled checks, per-state coverage, quarantine counts, recent diagnostics and the arrival log to the administration interface. Manual actions request a state history rebuild, update, census or local result refresh. The interface was developed by \auravname; the job state, diagnostics and data-management actions described here are the operational support beneath it.

Admin data and mutation routes require a local client address and reject non-local origins. This is a local-machine restriction, not a login or role system. The default service binds to \texttt{127.0.0.1}. A public reverse-proxy deployment requires deliberate authentication and proxy-aware access controls rather than assuming the current local address check remains sufficient.

\section{The daily operating cycle}
\begin{enumerate}
  \item Start the Python service against the intended output directory. It opens storage, recovers interrupted work, registers triggers and performs catch-up checks.
  \item Keep the process running with the computer awake and connected. The browser may be closed; it is not the scheduler.
  \item At 07:00 IST, the scheduler queues one update per state. Duplicate pending work is reused, and the worker starts eligible jobs as earlier work permits.
  \item The worker downloads recent windows, stores accepted rows, records failures and refreshes downstream results. Throttling can extend completion well beyond the trigger time.
  \item The backend exposes the committed observations and success markers; periodic website reads display the new state.
  \item Inspect the admin status when a run is paused, unusually slow or missing. A persistent source failure needs investigation before manually resuming collection.
\end{enumerate}

\section{What maintains the service, and what remains manual}
Automatic maintenance includes repeated collection, deduplication, retries, restart recovery, cache refresh and basic diagnostics. The process is not installed by this implementation as an operating-system service, and the database has no built-in scheduled off-machine backup. If the terminal session or computer stops the server, periodic callbacks stop too. Startup recovery only becomes active when the application is launched again.

An always-on deployment should add process supervision, controlled restarts, persistent storage, backups and external checks for a stopped service. This report describes these as future operational improvements. It does not equate a locally working scheduled application with an unattended public deployment.

\chapter{Data Coverage and Reporting Lag}

\section{Historical collection coverage}
The collection snapshot retained from 27 October 2026 contains approximately 685,000 validated observations across the two maintained states. Table~\ref{tab:data} records the original snapshot counts; it is not a live query of the continually updated database. The exact collected period varies with the state and build time, so the state-specific dates are retained.

\begin{table}[H]
\centering
\caption{Data collected (as of 27 October 2026)}
\label{tab:data}
\small
\begin{tabular}{@{}lrr@{}}
\toprule
& Uttar Pradesh & Rajasthan \\
\midrule
Validated price rows & about 488{,}000 & about 197{,}000 \\
Period & 27 Oct 2025 -- 27 Oct 2026 & 28 Oct 2025 -- 27 Oct 2026 \\
Commodities & 97 & 90 \\
Markets / districts & 251 / 73 & 160 / 37 \\
Series reported in the last 30 days & 4{,}912 & 1{,}732 \\
\bottomrule
\end{tabular}
\end{table}
A small Maharashtra pilot (about 8{,}700 rows) is kept in the database but is not shown.

\section{Reporting lag}
Table~\ref{tab:lag} and Figure~\ref{fig:lag} show the number of rows stored for each recent data date shortly after an update on 27 October 2026. Counts for Uttar Pradesh fall from about 1,300--1,600 per date up to 21 October to under 200 per date from 22 October, and keep shrinking towards the present: the newest days are not yet uploaded. Rajasthan's counts are much more complete for recent days. Sundays (18 and 25 October) have fewer rows in both states.

\begin{table}[H]
\centering
\caption{Stored rows per data date, 27 October 2026 snapshot}
\label{tab:lag}
\small
\begin{tabular}{@{}lrrlrr@{}}
\toprule
Date & UP & Rajasthan & Date & UP & Rajasthan \\
\midrule
13 Oct & 1{,}597 & 807 & 21 Oct & 1{,}442 & 879 \\
14 Oct & 1{,}614 & 798 & 22 Oct & 199 & 840 \\
15 Oct & 1{,}531 & 811 & 23 Oct & 198 & 723 \\
16 Oct & 1{,}568 & 744 & 24 Oct & 173 & 456 \\
17 Oct & 1{,}592 & 461 & 25 Oct (Sun) & 93 & 121 \\
18 Oct (Sun) & 1{,}266 & 108 & 26 Oct & 95 & 580 \\
19 Oct & 1{,}351 & 536 & 27 Oct & 30 & 64 \\
20 Oct & 1{,}304 & 461 & & & \\
\bottomrule
\end{tabular}
\end{table}

\begin{figure}[H]
\centering
\safefig{04_arrival_lag.png}{0.92}
\caption{Rows stored per data date for Uttar Pradesh and Rajasthan: recent dates in Uttar Pradesh are still largely missing.}
\label{fig:lag}
\end{figure}

The backend implication is that repeated collection must revisit older dates, not merely request today. A source can publish new information for an earlier market date even when no correction to an already stored price occurs. The database therefore preserves gaps and records the retrieval time separately from the market date.

\begin{figure}[H]
\centering
\safefig{05_example_series.png}{0.92}
\caption{An example stored price series (fixed market, commodity, variety and grade) with unreported days left as gaps, not filled.}
\label{fig:example}
\end{figure}


\section{The arrival log as operational evidence}
After a successful daily or weekly update, \texttt{app\_arrivals} records the number of stored observations per recent market date for that state, the newest date and the number of changed series. Counts cover dates from today minus 21 days onward. The administration query compares each snapshot with the preceding snapshot for the same state and displays positive row-count increases.

This distinguishes newly arriving reports for older dates from a simple advance in the latest date. It also makes timing decisions measurable: if a state's rows consistently arrive after the current check, a later check may be more useful. The count difference measures additional stored rows; it does not count every correction, because a changed value can leave the row count unchanged. The changed-series figure is therefore useful alongside it.

\section{Missing data and limits of a successful scrape}
Missing dates remain missing rather than being filled with zero, interpolation or a repeated price. The lack of a record does not identify whether the market was closed, did not trade that item or uploaded late. A complete download is complete relative to the requested source response, not a guarantee of exhaustive real-world market coverage.

The application normally lists series that have reported within 30 days. Older observations remain in the database, and hiding a series from the current picker does not delete its history. The backend supplies the reporting history and freshness information; eligibility for a statistical forecast is determined by the forecasting layer handled by \auravname.

\chapter{Collection Evidence and Backend Verification}

\section{Historical request budget and collection effort}
\label{sec:perf}
Table~\ref{tab:effort} retains measurements from the original collection phase. These illustrate the effect of state-wide windows and the sensitivity to throttling. They should not be read as a fixed daily completion-time guarantee.

\begin{table}[H]
\centering
\caption{Historical collection measurements (27 October 2026 evidence)}
\label{tab:effort}
\small
\begin{tabular}{@{}p{0.48\textwidth}p{0.20\textwidth}p{0.20\textwidth}@{}}
\toprule
Measurement & Uttar Pradesh & Rajasthan \\
\midrule
History commodity windows requested & 408 & 150 \\
History wall-clock time & About 12 hours & About 2.5 hours \\
Recent-date update, observed example & About 8 minutes, not throttled & About 52 minutes, throttled \\
Commodities in those update examples & 105 & 94 \\
\bottomrule
\end{tabular}
\end{table}

Across the two states, the historical record described about 1,100 completed windows: 602 with data, 17 with quarantined rows and 478 empty. Empty windows are a substantial cost, supporting active-commodity updates and gradual discovery. The same evidence recorded nine throttling and three CAPTCHA diagnostic events. A diagnostic event count is not a calibrated failure probability, but it identifies throttling as a practical bottleneck in those runs.

\section{Recorded daily operation on 6 November 2026}
The local SQLite job records were inspected read-only on 6 November 2026. Both state updates were created at 01:30 UTC, corresponding to the configured 07:00 IST trigger, and both finished successfully (Table~\ref{tab:dailyactual}). The times are presented in IST and rounded to the nearest displayed second.

\begin{table}[H]
\centering
\caption{Successful daily jobs recorded on 6 November 2026 (IST)}
\label{tab:dailyactual}
\small
\begin{tabular}{@{}lrrrr@{}}
\toprule
State & Queued & Completed & Commodities & Recorded failures \\
\midrule
Uttar Pradesh & 07:00:00 & 08:25:54 & 105 & 0 \\
Rajasthan & 07:00:00 & 09:45:10 & 94 & 1 \\
\bottomrule
\end{tabular}
\end{table}

The corresponding arrival records reported 19 changed Uttar Pradesh series and two changed Rajasthan series, with the newest stored market date of 6 November for each state. These are changed-series counts, not counts of new rows. Rajasthan's job recorded one failed attempt before success; the successful job row alone does not identify the earlier failure's cause. Its final attempt began at approximately 09:21:32 IST, so the creation and completion times also include queueing and retry effects.

These records establish that the backend completed the configured daily workflows on that date. They do not establish uninterrupted daily operation since installation: earlier records include a paused job and jobs that were executed well after creation. The report therefore distinguishes demonstrated runs from continuous-availability claims.

\section{Failure modes encountered and resolved}
\label{sec:failures}
\begin{enumerate}
  \item \textbf{Throttling mistaken for CAPTCHA rejection.} The source reused an error code, so the client had to classify on message content and stop the sweep instead of spending more challenge attempts.
  \item \textbf{Network outage treated as a generic exception.} A generic one-retry policy paused collection too early. Explicit outage classification permits deferred retries without changing the request pace.
  \item \textbf{Stale data used during downstream refresh.} A display-cache optimisation exposed an older list of series to processing. Fresh-read contexts now separate backend processing from temporarily stale display lists.
  \item \textbf{Cached results surviving logic changes.} Including model version in the result key only works if that version changes with the downstream logic. The integration contract now records that requirement.
  \item \textbf{Probes competing with routine collection.} Independent source investigation consumed the same request allowance as the running service. Investigative probes must be coordinated with scheduled collection.
\end{enumerate}

\section{Automated verification already present in the project}
The source tree contains 80 test methods across three test modules: 33 collector tests, 39 service tests and eight API tests. Some service tests concern the forecasting component handled by \auravname; the total describes the integrated repository and is not a claim of sole authorship for all tests. Relevant backend checks include:
\begin{itemize}
  \item supported OCR glyphs, rejected malformed images, bounded challenges and ticket scope;
  \item source parsing, explicit empty responses, units, invalid dates and request-dimension mismatches;
  \item idempotent ingestion, correction history, quarantine and transaction rollback;
  \item pagination integrity, deferred work, cooldowns, pace changes and interrupted-job recovery;
  \item active and census commodity selection, independent state settings and job deduplication;
  \item cache freshness, history export, local admin access and page reads that make no source requests.
\end{itemize}

The tests use fixtures, temporary databases and a fake source, allowing failures to be reproduced without consuming the portal's request allowance. They support behavioural verification but do not substitute for live throughput measurements or prove future source compatibility. Reproduction commands are given in Appendix~D; the existing suite is described here without claiming a fresh execution for this document revision.

\chapter{Limitations and Future Work}

\section{Current limitations}
\begin{itemize}
  \item \textbf{Single-machine availability.} The scheduler lives inside the running server; sleep, shutdown or process termination stops callbacks. There is no installed process supervisor or external uptime monitor.
  \item \textbf{Source-controlled freshness.} Successful collection cannot retrieve observations that the market has not yet uploaded. A recent maximum date is not evidence that every series is current.
  \item \textbf{Non-preemptive work.} A long running job can delay a daily update or local rollover. The queue does not interrupt work or reserve a fast execution lane for local-only refreshes.
  \item \textbf{Bounded recovery.} Persistent checkpoints survive restarts, but partial pages are re-downloaded. Retries crossing a date boundary can create new interval keys. Missed weekly triggers are not reconstructed automatically.
  \item \textbf{Global pause.} A fatal state-specific error pauses scheduled collection for all tracked states. An isolated per-state pause would improve independence.
  \item \textbf{Catalogue and format dependence.} Cached filters need explicit refresh, and changes to response fields, CAPTCHA rendering or limits may require code changes.
  \item \textbf{Storage growth and backups.} Raw-update archiving is limited, but job history, diagnostics and result caches still accumulate. Automatic retention and backup/restore testing are not implemented.
  \item \textbf{Deployment boundary.} Local client-address restrictions do not provide production user authentication. Public hosting needs a separate access and process-management setup.
\end{itemize}

\section{Backend improvements}
\begin{enumerate}
  \item \textbf{Always-on operation.} Deploy with persistent storage, process supervision, deliberate shutdown handling, HTTPS, authenticated administration and external freshness checks. Verify source access from the chosen host before relying on it.
  \item \textbf{Backups and retention.} Add scheduled SQLite backups, archive retention and restoration exercises. Preserve enough provenance to reproduce reported collection evidence.
  \item \textbf{Better freshness scheduling.} Use per-state arrival logs to choose check times. Consider a second check only where new observations justify the additional request cost.
  \item \textbf{Stronger job isolation.} Separate local-only processing from long source downloads, introduce per-state pause control and persist last-completed weekly markers for deliberate catch-up.
  \item \textbf{Catalogue maintenance.} Refresh and version the source filter catalogue, detect removals or renamed entries and review changes before collection continues.
  \item \textbf{Broader coverage.} Test an all-states collection strategy. Historical probes showed an all-states filter; estimated history collection could fall from about 6,000 state-wise reports to 1,500--2,000 broader reports. These are planning estimates, not measured nationwide operation.
  \item \textbf{Longer histories.} Earlier-year wheat probes demonstrated available rows and substantial market-name continuity. Collecting those years requires the same validation and provenance controls; any benefit to forecasting must be evaluated by the forecasting contributor.
\end{enumerate}

\chapter{Conclusion}
The project phases documented in this report produced a backend that turns individually requested Agmarknet reports into a validated, queryable and repeatedly refreshed data service. Source investigation established that report count and throttling are the main collection constraints. State-wide commodity windows, controlled ticket reuse and pagination checks reduce unnecessary access while preserving the source's dimensions.

The database supplies consistent observations, quarantine, correction history and provenance. Persistent jobs, exact-window checkpoints, adaptive pacing and retry classification make collection recoverable. The in-process scheduler queues daily and weekly work, while the worker serialises execution and the data-serving layer keeps existing history available during updates. Recorded successful jobs on 6 November 2026 demonstrate this daily workflow without implying continuous operation or complete source reporting.

The forecasting and website-development phases, including model choice and predictive evaluation, were handled by \auravname. This report has confined its technical claims to the data pipeline and backend operation supporting those components. The next backend steps are supervised always-on deployment, tested backups, better per-state scheduling and controlled expansion of collection coverage.

\appendix
\chapter{Database Schema and Data Lifecycle}
\begin{lstlisting}[caption={Simplified schema; fields omitted where not central}]
observations(key PK, date, state, district, market, commodity,
             variety, grade, min_price, max_price, modal_price,
             unit, payload)
series_state(id PK, dimensions, revision, updated)
batches(id PK, state_id, commodity_id, start, end, market_id,
        parent, status, attempts, rows, valid, rejected, error, updated)
pages(batch_id, run_id, page, path, sha256, rows, total,
      PK(batch_id, run_id, page))
quarantine(id PK, batch_id, reason, raw, created)
revisions(id PK, observation_key, previous_payload, new_payload, created)
events(id PK, created, batch_id, event, detail)
app_runs(id PK, kind, state, status, attempt, created, started,
         finished, retry_at, detail, progress, total)
app_windows(run_id, window_key, status, batch_id,
            PK(run_id, window_key))
app_settings(key PK, value)
app_diagnostics(id PK, run_id, created, payload)
app_arrivals(id PK, run_id, state, kind, at, changed_series,
             newest, counts)
app_forecasts(series_id, target_date, revision, model_version,
              generated, payload,
              PK(series_id, target_date, revision, model_version))
\end{lstlisting}

The payload of a canonical observation retains its source metadata. \texttt{app\_forecasts} is a downstream integration cache; its presence does not assign authorship of the forecasting component to this report. Durable database rows and disk archives survive process restarts; in-memory list caches and report tickets do not.

\chapter{Backend API and Responsibility Boundary}
\begin{longtable}{@{}p{0.44\textwidth}p{0.48\textwidth}@{}}
\toprule
Endpoint & Backend purpose \\
\midrule
\endfirsthead
\toprule
Endpoint & Backend purpose \\
\midrule
\endhead
\texttt{GET /api/status} & Per-state completion markers, date, collection pause and current/waiting work. \\
\texttt{GET /api/selections?state=} & Saved series dimensions for one state and availability flags supplied by downstream integration. \\
\texttt{GET /api/series/\{id\}/history?days=} & Date-ordered stored observations. \\
\texttt{GET /api/series/\{id\}/history.csv} & Export of a series' complete saved history and dimensions. \\
\texttt{GET /api/series/\{id\}/forecast} & Exact-date/revision downstream payload; queue a local refresh if missing. Statistical content is supplied by the forecasting phase handled by \auravname. \\
\texttt{GET /api/admin/status} & Local operational status, diagnostics, coverage, arrivals and scheduled next runs. \\
\texttt{POST /api/admin/enqueue?kind=\&state=} & Local manual collection or refresh action. Bootstrap requires a chosen state. \\
\bottomrule
\end{longtable}

The root page, admin page and static assets are the integrated website interface. They are included in the running application but their presentation and front-end development were handled by \auravname. API reads are not a source-collection trigger.

\chapter{Startup and Retry Algorithms}
\begin{lstlisting}[caption={Startup recovery and registration, simplified}]
if collector file lock is available:
    change interrupted running jobs back to queued
register daily_update at 07:00 IST
register weekly_census at Sunday 03:00 IST
register local_result_rollover at 00:05 IST
warm saved-data lists in a background thread
start scheduler
for each tracked state:
    if initial history incomplete: enqueue bootstrap
    else if last success absent or older than 24 hours: enqueue update
enqueue one shared local refresh
start single worker

worker:
    while service is running:
        drain all eligible jobs using the selection policy
        wait up to 10 seconds or until shutdown signal
\end{lstlisting}

\begin{lstlisting}[caption={Failure transition after a running job, simplified}]
if access denied or source format changed:
    status = paused; collection_paused = true
else if deferred / rate limited / outage:
    if prior_failures < 12:
        delay = exception_base_minutes * 2**min(prior_failures, 1)
        status = retry_wait; retry_at = now + delay
    else:
        status = paused; collection_paused = true
else:
    retry once after 30 minutes; otherwise pause collection
increment failure counter; save detail and timestamps
keep committed observations and completed exact-window checkpoints
\end{lstlisting}

\chapter{Reproducibility and Evidence Dates}
\section*{Starting the local backend}
The application expects Python 3.11 or later on macOS or Linux. Dependencies for the integrated service are listed in \path{requirements-app.txt}. From the project directory:
\begin{lstlisting}
cd agmarknet_collector
.venv/bin/python -m pip install -r requirements-app.txt
.venv/bin/python -m agmarknet serve --output output
# local service: http://127.0.0.1:8000
# local admin:   http://127.0.0.1:8000/admin
\end{lstlisting}
The environment must already exist; if needed it can be created with \texttt{python3 -m venv .venv}. Starting the server can queue real collection through startup catch-up, so use the existing intended output directory and avoid running multiple unmanaged instances.

\section*{Inspecting and testing}
\begin{lstlisting}
# Deterministic offline verification using temporary databases:
.venv/bin/python -m unittest discover -s tests -v

# Read-only inspection of recent persisted jobs, without starting the service:
sqlite3 -readonly output/agmarknet.sqlite \
  "SELECT kind,state,status,created,finished,attempt FROM app_runs ORDER BY created DESC LIMIT 10;"
\end{lstlisting}
UTC timestamps in SQLite require conversion by adding 05:30 for the IST examples in Chapter~7. The read-only inspection above does not enqueue jobs or refresh results.

\section*{Source files for the documented behaviour}
\begin{itemize}
  \item \path{agmarknet/client.py} and \path{agmarknet/ocr.py}: source requests, ticket scope, challenge handling and request-level retries.
  \item \path{agmarknet/storage.py}: normalisation, identity, transactions, quarantine, revisions and export.
  \item \path{agmarknet/service.py}: application jobs, report windows, census selection, pace, persisted state and cache coordination.
  \item \path{agmarknet/app.py}: runtime, scheduler registration, catch-up, data APIs and local administration boundary.
  \item \path{agmarknet/__main__.py}: command-line entry points and server startup.
\end{itemize}

Historical coverage, request-budget, collection-duration and reporting-lag tables are retained from the original 27 October 2026 evidence. The daily-run table was read from the local database on 6 November 2026. The implementation descriptions were checked against the source during this revision. Historical figures should be regenerated before replacing their underlying snapshot counts with current totals; they must not be relabelled as live measurements.

\begin{thebibliography}{99}
\addcontentsline{toc}{chapter}{Bibliography}

\bibitem{agmarknet}
Government of India, Ministry of Agriculture and Farmers Welfare,
\emph{Agmarknet: Agricultural Marketing Information Network (Agmarknet 2.0)}.
\url{https://agmarknet.gov.in/}.

\bibitem{glez2014}
D. Glez-Pe\~na, A. Louren\c{c}o, H. L\'opez-Fern\'andez, M. Reboiro-Jato and F. Fdez-Riverola,
\emph{Web scraping technologies in an API world},
Briefings in Bioinformatics, vol. 15, no. 5, pp. 788--797, 2014.

\bibitem{sqlitewal}
SQLite Documentation, \emph{Write-Ahead Logging}.
\url{https://www.sqlite.org/wal.html}.

\bibitem{sqlitebackup}
SQLite Documentation, \emph{SQLite Backup API}.
\url{https://www.sqlite.org/backup.html}.

\bibitem{apscheduler}
APScheduler 3.x Documentation, \emph{User guide: missed job executions, coalescing and concurrent instances}.
\url{https://apscheduler.readthedocs.io/en/3.x/userguide.html}.

\bibitem{fastapi}
S. Ram\'irez, \emph{FastAPI documentation}.
\url{https://fastapi.tiangolo.com/}.

\end{thebibliography}
\end{document}
